In this section, we introduce the main tools that we will use for both linear-algebraic and combinatorial algorithms in later sections.
Until \autoref{subsec:non-sc}, we fix an unambiguous, complete, and strongly connected monoid morphism~$M$. In \autoref{subsec:non-sc} we will show that the case where $M$ is not strongly connected can be easily reduced to the strongly connected case.


\subsection{Columns, rows and the structure of minimum rank matrices}\label{subsec:col-row-defs}

The concept of maximum columns and rows plays a crucial role in dealing with reachability problems in unambiguous monoid morphisms. Abusing language slightly in the following, by \emph{column} we refer to column vectors of the form $[w \cdot q] = M(w)[q] \in \{0,1\}^Q$ where $w \in \Sigma^*$ and $q \in Q$.
Similarly, a \emph{row} is of the form $[q \cdot w]^T = [q]^T M(w)$. See \autoref{fig:ex-columns-rows} for an example. In the case of total DFAs, all rows are of the form $[q]^T$. This fact makes total DFAs significantly simpler to deal with than general complete UFAs.

\begin{figure}[h]
\centering
% Center the two drawings as a group
\begin{tabular}{@{}c@{\hspace{1cm}}c@{}}
$
M(a) = \begin{pmatrix} 1 & 0 & 1 & 0 \\ 1 & 0 & 1 & 0 \\ 0 & 0 & 0 & 0 \\ 0 & 0 & 0 & 0 \end{pmatrix} \quad 
M(b) = \begin{pmatrix} 0 & 0 & 0 & 0 \\ 0 & 0 & 0 & 0 \\ 0 & 1 & 0 & 1 \\ 0 & 1 & 0 & 1 \end{pmatrix}
$
&
\begin{tikzpicture} [baseline=(current bounding box.center), node distance = 2cm]
%,
\tikzset{every state/.style={inner sep=1pt,minimum size=1.5em}}

\node [state] at (-2, 0) (1) {$1$};
\node [state] at (0, 1) (2) {$2$};
\node [state] at (2, 0) (3) {$3$};
\node [state] at (0, -1) (4) {$4$};
\path [thick,->,>=stealth]

(2) edge [] node[above] {$a$} (1)
(2) edge [bend left=20] node[above] {$a$} (3)
(1) edge [] node[above, pos=0.3] {$a$} (3)

(1) edge [loop below] node[left] {$a$} (1)

(3) edge [bend left=20] node[above] {$b$} (2)
(3) edge [] node[below] {$b$} (4)
(4) edge [] node[left, pos=0.3] {$b$} (2)

(4) edge [loop,out=200,in=160,looseness=8] node[left] {$b$} (4)
;
\end{tikzpicture}
\end{tabular}
\caption{$[a \cdot 3] = M(a) [3] = [\{1, 2\}]$ is a column; $[2 \cdot a]^T = [2]^T M(a) = [\{1, 3\}]^T$ is a row.}\label{fig:ex-columns-rows}
\end{figure}

% Fix indentation
\noindent
A column $[C]$ is called \emph{maximal} if there is no column $[C']$ such that $[C'] \ne [C]$ and $[C'] \ge [C]$ (that is, $C \subset C'$). Maximal rows are defined in the same way. Recall that the inequalities are taken component-wise. 

Let $A \in \{0,1\}^{m \times n}$ be a zero-one matrix.
One can view~$\rkr(A)$ as the least number~$r$ such that there are matrices $C \in \mathbb{R}^{m \times r}$ and $R \in \mathbb{R}^{r \times n}$ with $A = C R$.
Define the \emph{unambiguous rank} $\rku(A)$ as the least number~$r$ such that there are matrices $C \in \{0,1\}^{m \times r}$ and $R \in \{0,1\}^{r \times n}$ such that $A = C R$.
Analogously to $\rkr(M)$, define also $\rku(M) := \min\{\rku(M(w)) \mid w \in \Sigma^*\}$.
Clearly, $\rkr(A) \le \rku(A)$, and the inequality can be strict, but in \autoref{cor:ranks-equal} below we show that $\rkr(M) = \rku(M)$.
The reason we are interested in the unambiguous rank is that \autoref{thm:Cesari} below implies that there is always a matrix with a very simple structure such that its unambiguous rank is equal to its real rank and both ranks are minimum.

In the following let us write $r := \rku(M)$ when $M$~is understood.
A word $u \in \Sigma^*$ is \emph{of minimum unambiguous rank} if $\rku(M(u)) = r$.
If $u \in \Sigma^*$ is of minimum unambiguous rank then so is $v u w$ for all $v, w \in \Sigma^*$. 

Words of unambiguous rank one, known as synchronising words, play an especially important role due to their applications in the theory of codes, as explained in \autoref{sec-existing}. It is easy to see that a word $w$ has unambiguous rank one if and only if there exist $C, R \subseteq Q$ such that $w$ maps a state $p$ to a state $q$ if and only if $p \in C$ and $q \in R$.
For total DFAs, we moreover have that $C = Q$ and $R$ has cardinality one.

%Changed to thmC
\begin{thmC}[C{\'{e}}sari \cite{Cesari74}] \label{thm:Cesari}
Let $u \in \Sigma^*$ be of minimum unambiguous rank.
There are pairwise disjoint sets $C_1, \ldots, C_r \subseteq Q$ and pairwise disjoint sets $R_1, \ldots, R_r \subseteq Q$ such that \[M(u) = \sum_{i=1}^r [C_i] [R_i]^T.\]
Moreover, each $[C_i]$ and $[R_i]^T$ is, respectively, a maximal column and a maximal row.
\end{thmC}

% Fix indentation
\noindent
This theorem will play a central role.
A proof can be found in \cite[Proposition 4]{BealCKP08}. For the sake of completeness, we provide an elementary proof of the next two sections.
In the case of a total DFA, $R_1$ is a singleton and \autoref{thm:Cesari} is fairly obvious.



In \autoref{thm:Cesari}, since the $C_i$ are pairwise disjoint and the $R_i$ are pairwise disjoint, each $[C_i] [R_i]^T$ forms, intuitively, a ``combinatorial rectangle'', and no such rectangle shares a row or a column with any other rectangle.
The column vectors~$[C_i]$ are exactly the nonzero columns of~$M(u)$ and linearly independent, and the row vectors~$[R_i]^T$ are exactly the nonzero rows of~$M(u)$ and linearly independent.
Thus, $r$~is the number of distinct nonzero columns and also the number of distinct nonzero rows in~$M(u)$.
It follows that $r = \rku(M(u)) = \rkr(M(u))$. Thus we have:

\begin{cor} \label{cor:ranks-equal}
We have $r = \rku(M) = \rkr(M)$.
\end{cor}

% Fix indentation
\noindent
We can thus define $\rkg(M)$ as $\rku(M) = \rkr(M)$.
For words $w \in \Sigma^*$ that are not of minimum unambiguous rank, we may have $\rkr(M(w)) < \rku(M(w))$, but the rank of such matrices will rarely play a role in the following. In what follows, we call words of minimum unambiguous rank simply words of minimum rank. Since below we never refer to the real rank of words, this will not lead to any confusion.

\subsection{A proof of the first statement in C{\'{e}}sari's theorem}

We use the following lemma.
\begin{lem} \label{lem:Cesari-lemma}
Let $A \in \{0, 1, 2, \ldots\}^{Q \times Q}$ be with $[Q]^T A \ge [Q]^T$ and $A [Q] \ge [Q]$.
Then $A$~is a permutation matrix \emph{or} $\sup_{i \ge 0} [Q]^T A^i [Q] = \infty$.
\end{lem}
\begin{proof}
By a straightforward induction argument we have that $A^{i+1} [Q] \ge A^i [Q]$ holds for all $i \ge 0$.
It suffices to show that $A$~is a permutation matrix \emph{or} for all $i \ge 0$ we have $[Q]^T A^{i+1} [Q] > [Q]^T A^i [Q]$.
We suppose that $A$~is not a permutation matrix and show by induction on~$i$ that $[Q]^T A^{i+1} [Q] > [Q]^T A^i [Q]$ holds for all $i \ge 0$.

Concerning the induction base, $i=0$, since $A$ is not a permutation matrix, we cannot have both $[Q]^T A = [Q]^T$ and $A [Q] = [Q]$ (in fact, we can have neither).
Therefore, $[Q]^T A [Q] > [Q]^T [Q]$.
For the induction step, suppose that $[Q]^T A^{i+1} [Q] > [Q]^T A^i [Q]$ holds for some $i \ge 0$.
Then $A^{i+1} [Q] > A^i [Q]$; i.e., the inequality $A^{i+1} [Q] \ge A^i [Q]$ is strict in some component.
Since $[Q]^T A \ge [Q]^T$, the vector $[Q]^T A$ is strictly positive in all components.
It follows that $[Q]^T A A^{i+1} [Q] > [Q]^T A A^i [Q]$, as required.
\end{proof}

% Fix indentation
\noindent
This allows us to prove the first statement of \autoref{thm:Cesari}. The second statement is proved at the end of the next subsection.

\begin{prop} \label{prop:Cesari-first}
Let $u \in \Sigma^*$ be of minimum rank.
There are pairwise disjoint sets $C_1, \ldots, C_r \subseteq Q$ and pairwise disjoint sets $R_1, \ldots, R_r \subseteq Q$ such that $M(u) = \sum_{i=1}^r [C_i] [R_i]^T$.
\end{prop}
\begin{proof}
Since $\rku(M(u)) = r$, there are $C \in \{0,1\}^{Q \times r}$ and $R \in \{0,1\}^{r \times Q}$ with $M(u) = C R$.
For notational convenience, take a finite set $X$ with $|X| = r$ and write $C \in \{0,1\}^{Q \times X}$ and $R \in \{0,1\}^{X \times Q}$.
Let $w \in \Sigma^*$.
Define $P(w) := R M(w) C \in \{0,1,2,\ldots\}^{X \times X}$.
For all $i \ge 0$ we have $C P(w)^i R = (C R M(w))^i C R = M((u w)^i u) \in \{0,1\}^{Q \times Q}$.
Since $C$~has no zero columns and $R$~has no zero rows, we have $P(w)^i \in \{0,1\}^{X \times X}$ for all~$i$.
Further, since $\rku(P(w)) \ge \rku(C P(w) R) = \rku(M(u w u)) = r$, we have $\rku(P(w)) = r$.
In particular, $P(w)$ does not have a zero row or column; i.e., $P(w) [X] \ge [X]$ and $[X]^T P(w) \ge [X]^T$. 
It follows from \autoref{lem:Cesari-lemma} that $P(w)$~is a permutation matrix.
As $w \in \Sigma^*$ was arbitrary, $P(w)$ is a permutation matrix for all~$w$.

Let $q_1 \in Q$, and let $q_2 \in Q$ and $y \in X$ be such that $[q_2]^T C [y] = 1$.
By strong connectedness, there is $w \in \Sigma^*$ with $[q_1]^T M(w) [q_2] = 1$.
Thus,
\[
P(w) [y] \ = \ R M(w) C [y] \ \ge \ R [q_1] [q_1]^T M(w) [q_2] [q_2]^T C [y] \ = \ R [q_1] \,.
\]
As $P(w)$ is a permutation matrix and $q_1 \in Q$ was arbitrary, it follows that all nonzero columns of~$R$ are of the form $[x]$ for $x \in X$.
Similarly, all nonzero rows of~$C$ are of the form~$[x]^T$ for $x \in X$.

For each $x \in X$, define $C_x, R_x \subseteq Q$ such that $[C_x] = C [x]$ and $[R_x]^T = [x]^T R$.
The~$R_x$ are pairwise disjoint, as if there was $q \in R_x \cap R_y$ with $x \ne y$, then $R [q] \ge [\{x,y\}]$, contradicting what we proved in the previous paragraph.
Similarly, the $C_x$ are pairwise disjoint.
Finally, we have
\[
 M(u) \ = \ C R \ = \ \sum_{x \in X} C [x] [x]^T R \ = \ \sum_{x \in X} [C_x] [R_x]^T\,,
\]
as desired.
\end{proof}


\subsection{The weight of columns and rows}\label{subsec:weights}

The results in this subsection, about the column and row vectors that appear in the matrices~$M(w)$, are mostly due to \cite{Cesari74}; see also \cite[Section 3]{BealCKP08}.
Since a notion of column and row weight will be crucial for us in the later development, we phrase and prove the results around these concepts, but we do not view the lemmas of this subsection as novel.

Define $\overline{A} = \frac{1}{|\Sigma|} \sum_{a\in \Sigma} M(a) \in [0,1]^{Q \times Q}$.
Since $M$ is strongly connected, $\overline{A}$~is irreducible.
Since $M$ is unambiguous, the spectral radius of~$\overline{A}$ is at most~$1$, and since $M$ is complete, it is at least~$1$.
Thus, the spectral radius of~$\overline{A}$ equals~$1$.
Since $\overline{A}$~is irreducible, it follows from basic Perron-Frobenius theory that $\overline{A}$~has an eigenvalue~$1$ and every right eigenvector with eigenvalue~$1$ is a multiple of a strictly positive vector, say $\beta \in \RP^Q$.
Since $\overline{A}$~has only rational entries, we can assume $\beta \in \mathbb{Q}_{>0}^Q$.
Similarly for left eigenvectors.
Therefore, there are $\alpha, \beta \in \mathbb{Q}_{>0}^Q$ with $\alpha^T \overline{A} = \alpha^T$ and $\overline{A} \beta = \beta$.
Without loss of generality, we assume that~$\alpha^T \beta = 1$.

In the total DFA case, since $M(a) [Q] = [Q]$ for all $a \in \Sigma$, we have $\overline{A} [Q] = [Q]$ and so it is natural to take $\beta = [Q]$.
In that case, $\alpha^T [Q] = \alpha^T \beta = 1$ means that $\alpha^T = \alpha^T \overline{A}$ is the (unique) stationary distribution of the Markov chain whose transition probabilities are given by the row-stochastic matrix~$\overline{A}$; intuitively, in this Markov chain a letter $a \in \Sigma$ is picked uniformly at random in every step.


Define the \emph{weight} of a column~$y$ and of a row~$x^T$ by $\alpha^T y \in \mathbb{R}$ and $x^T \beta \in \mathbb{R}$, respectively.
Denote the maximum column weight and the maximum row weight by $\mcw$ and $\mrw$, respectively, i.e.,
\[
 \mcw \ := \ \max\{\alpha^T y \mid y \text{ is a column}\} \quad \text{ and } \quad
 \mrw \ := \ \max\{x^T \beta   \mid x^T \text{ is a row}\}\,.
\]

% Fix indentation
\noindent
A column~$y$ is called \emph{of maximum weight} if $\alpha^T y = \mcw$, and analogously for rows.
In the total DFA case, every row is of the form $[q]^T$ for some $q \in Q$, hence every row is of maximum weight.


For $q \in Q$ define
\[
 \Mer(q) \ := \ \{q' \in Q \mid \exists\,S \supseteq \{q,q'\} \text{ such that } [S] \text{ is a column}\}\,.
\]
Intuitively, $\Mer(q)$ consists of the states that can ``appear'' in a column together with~$q$, or, equivalently, the states that are ``mergeable'' with~$q$ (that is, can be mapped to the same state by a word).
Note that $q \in \Mer(q)$. 

The following lemmas apply symmetrically also to rows.

\begin{lem} \label{lem:max-column}
Let $v \in \Sigma^*$ and $q \in Q$ be such that $[v \cdot q]$ is a column of maximum weight. Then we have:
\begin{enumerate}
\item[(a)] $[u v \cdot q]$ is a column of maximum weight for all $u \in \Sigma^*$;
\item[(b)] $v \cdot q' = \emptyset$ holds for all $q' \in \Mer(q) \setminus \{q\}$.
\end{enumerate}
\end{lem}

\begin{proof}
Towards~(a), let $a \in \Sigma$.
It suffices to prove that $M(a) [v \cdot q]$ is of maximum weight.
We have $\alpha^T [v \cdot q] = \alpha^T \overline{M} [v \cdot q] = \frac{1}{|\Sigma|} \sum_{b \in \Sigma} \alpha^T M(b)[v \cdot q]$.
Thus, if $\alpha^T M(a) [v \cdot q] < \alpha^T [v \cdot q]$, then there would also exist $b \in \Sigma$ with $\alpha^T M(b) [v \cdot q] > \alpha^T [v \cdot q]$, contradicting that $[v \cdot q]$ is of maximum weight.
So $M(a)[v \cdot q]$ is of maximum weight.

Towards~(b), let $q' \in \Mer(q) \setminus \{q\}$.
Since $M$ is strongly connected, there is $w \in \Sigma^*$ with $[w \cdot q] \ge [q] + [q']$.
Thus, $[v w \cdot q] \ge [v \cdot q] + [v \cdot q']$.
It follows that, since $[v \cdot q]$ is of maximum weight, so is $[v w \cdot q]$.
Hence, $[v \cdot q'] = 0$.
\end{proof}

\begin{lem} \label{lem:non-max-column}
Let $S \subseteq Q$ be such that $[S]$ is a column which is not of maximum weight.
\begin{enumerate}
\item[(a)] There is $S' \supsetneq S$ such that $[S']$ is a column of maximum weight.
\item[(b)] There is $u \in \Sigma^*$ with $M(u) [S] = 0$.
\end{enumerate}
\end{lem}


\begin{proof}
Let $v \in \Sigma^*$ and $q \in Q$ be such that $[v \cdot q]$ is not of maximum weight.

Towards~(a), let $w \in \Sigma^*$ and $t \in Q$ be such that $[w \cdot t]$ is of maximum weight and $w \cdot t \ni q$.
By \autoref{lem:max-column}~(a), $[v w \cdot t]$ is of maximum weight.
Since $v w \cdot t \supseteq v \cdot q$ and $[v \cdot q]$ is not of maximum weight, we have $v w \cdot t \supsetneq v \cdot q$.

Towards~(b), let $p \in S' \setminus S$.
We have $S \subseteq S' \subseteq \Mer(p)$ and thus $S \subseteq \Mer(p) \setminus \{p\}$.
Let~$u \in \Sigma^*$ be such that $[u \cdot p]$ is of maximum weight.
Then it follows from \autoref{lem:max-column}~(b) that $M(u) [S] = 0$.
\end{proof}

\begin{lem} \label{lem:max-notions}
    A column (respectively, row) is maximal if and only if it is of maximum weight.
\end{lem}
\begin{proof}
If a column is of maximum weight, it is clearly maximal.
Conversely, if a column is not of maximum weight, then \autoref{lem:non-max-column}~(a) says it is not maximal.
\end{proof}

% Fix indentation
\noindent
An important property that we will need later is that the set of maximal columns is closed under left multiplication by matrices from the monoid, as stated in the following lemma. Note that this is no longer true without the completeness assumption, and is the key reason why the case of complete matrix monoids is easier to deal with.

\begin{lem} \label{lem:max-column-stable}
Let $v \in \Sigma^*$ and $q \in Q$ be such that $[v \cdot q]$ is a maximal column. Then $[u v \cdot q]$ is a maximal column for all $u \in \Sigma^*$.
\end{lem}
\begin{proof}
Immediate from \autoref{lem:max-column}~(a) and \autoref{lem:max-notions}.
\end{proof}

% Fix indentation
\noindent
The following lemma will be useful later to construct minimum rank matrices from maximal columns and rows.

\begin{lem} \label{lem:all-max-column}
Let $w \in \Sigma^*$ be such that all nonzero columns and rows in $M(w)$ are maximal. Then %$M(ww)$
$w w$ is of minimum rank.
\end{lem}
\begin{proof}
Define $S := \{q \in Q \mid w \cdot q \ne \emptyset \ne q \cdot w\}$.
Since
\[
M(w w) \ = \ M(w) I M(w) \ = \ M(w) \big( \sum_{q \in Q} [q] [q]^T \big) M(w) \ = \ \sum_{q \in Q} [w \cdot q] [q \cdot w]^T \ = \ \sum_{q \in S} [w \cdot q] [q \cdot w]^T\,,
\]
we have $\rku M(w w) \le |S|$. %$ \le \rku M(u w w)$.

% Fix indentation
\noindent
Let $u \in \Sigma^*$ be of minimum rank.
Suppose there were a state $p \in Q$, two states $q, q' \in S$ and $u w$-labelled paths from $p$ to~$q$ and also from $p$ to~$q'$.
Then $p \cdot u w w \supseteq q \cdot w \cup q' \cdot w$, contradicting the maximality of the row $[q \cdot w]^T$.
Therefore, the sets $u w \cdot q$, where $q \in S$, are pairwise disjoint.
Since the columns $[w \cdot q]$ for $q \in S$ are maximal, by \autoref{lem:max-column-stable} the columns $[u w \cdot q]$ are also maximal and in particular nonzero.
Moreover, these columns appear in $M(u w w)$, as $q \cdot w \ne \emptyset$ for all $q \in S$.
Therefore, $|S| \le \rku M(u w w)$.
Since $u$ is of minimum rank, we have $\rku M(u w w) \le r$.
By combining all inequalities we obtain that $\rku M(w w) \le r$.
Hence, $w w$ is of minimum rank.
\end{proof}

% Fix indentation
\noindent
We can now complete the proof of \autoref{thm:Cesari}.
\begin{proof}[Proof of \autoref{thm:Cesari}]
The first statement is \autoref{prop:Cesari-first}.
We prove the second statement only for columns, as the proof for rows is analogous.
Towards a contradiction, suppose that~$[C_i]$ is a non-maximal column; without loss of generality, say $i=r$.
By \autoref{lem:max-notions}, $[C_i]$ is not of maximum weight.
By \autoref{lem:non-max-column}~(b), there is $v \in \Sigma^*$ such that $M(v) [C_r] = 0$.
Then 
\[
 M(v u) \ = \ \sum_{i=1}^r M(v) [C_i] [R_i]^T \ = \ \sum_{i=1}^{r-1} M(v) [C_i] [R_i]^T \ = \ C R^T\,,
\]
where 
$C, R \in \{0,1\}^{Q \times (r-1)}$ are the matrices whose $i$th columns are $M(v) [C_i]$ and $[R_i]$, respectively.
Hence, $\rku(M(v u)) \le r-1$, contradicting the definition of~$r$.
\end{proof}


\subsection{Weight preservation property and minimum rank}\label{subsec:weight-preserv}

Every word $w$ of minimum rank in a total DFA induces a partition of the state set into subsets of states mapped by $w$ to the same state (that is, into columns). It was observed by Friedman in \cite{Friedman1990} that all sets of such a partition have the same weight. This observation has many applications to variations of the road colouring problem \cite{Friedman1990,Kari2001,Gusev2016,Kari2019}. Moreover, it was proved in \cite[Theorem 6]{Gusev2016}, again in connection with road colouring, that for every $w$ the weights of all columns in $M(w)$ sum up to $1$ (assuming $\beta = [Q]$ as suggested previously). 
This can be seen as a weight preservation property: the total weight of columns in the matrix of a word is preserved under multiplication by any matrix from the monoid.
As a result we get that $1 = r \cdot \mcw$, and hence~$r = \frac{1}{\mcw}$. The proof of the weight preservation property for total DFAs is quite simple and relies on the fact that for a state $q$ and a word~$w$ the set~$q \cdot w$ is always a singleton. For complete UFAs this is no longer true; in particular, $q \cdot w$ can be the empty set, thus permanently ``losing'' some weight collected in $q$. Hence, a more refined property is required. The following result provides such a property.
It also turns out that its proof requires more sophisticated techniques than in the total DFA case.


\begin{thm} \label{thm:rank-weight}
For all $w \in \Sigma^*$ we have $1 = \alpha^T M(w) \beta = r \cdot \mcw \cdot \mrw$.
\end{thm}

% Fix indentation
\noindent
Similarly to the total DFA case, this result allows us to reduce computing~$r$ to computing $\mcw$ and~$\mrw$, which we will use later in our algorithms. Recall that we have defined $\alpha^T$ and~$\beta$ so that $\alpha^T \beta = 1$.



Towards a proof of \autoref{thm:rank-weight} we first prove the following lemma.

\begin{lem} \label{lem:minrank-weight}
Let $u \in \Sigma^*$ be of minimum rank.
Then $\alpha^T M(u) \beta = r \cdot \mcw \cdot \mrw$.
\end{lem}
\begin{proof}
Let $M(u) = \sum_{i=1}^r [C_i] [R_i]^T$ be as in \autoref{thm:Cesari}.
Each $[C_i]$ and each $[R_i]$ is of maximum weight.
Thus,
\[
 \alpha^T M(u) \beta \ = \ \sum_{i=1}^r \alpha^T [C_i] [R_i]^T \beta \ = \ \sum_{i=1}^r \mcw \cdot \mrw \ = \ r \cdot \mcw \cdot \mrw\,. \qedhere
\]
\end{proof}

% Fix indentation
\noindent
We also need the following proposition.

\begin{prop}\label{prop:compact}
Let $x \in \mathbb{R}^{Q}$ and $c \in \mathbb{R}$ be such that $x^T M(u) \beta = c$ holds for all $u \in \Sigma^*$ of minimum rank.
Then $x^T M(w) \beta = c$ holds for all $w \in \Sigma^*$.
\end{prop}

\begin{proof}
Let $w \in \Sigma^*$.
Let $u \in \Sigma^*$ be of minimum rank.
Recall that every word that contains~$u$ as a factor is of minimum rank.
For every $k \ge 0$, partition $\Sigma^k$ into sets $W_0(k)$ and $W_1(k)$ so that $W_0(k) = \Sigma^k \cap (\Sigma^* u \Sigma^*)$ and $W_1(k) = \Sigma^k \setminus (\Sigma^* u \Sigma^*)$; i.e., $W_0(k), W_1(k)$ are the sets of length-$k$ words that do or do not contain~$u$ as a factor, respectively.
For all $v \in W_0(k)$ both~$v$ and $w v$ are of minimum rank.
Thus, we have $x^T M(w v) \beta = c$ for all $v \in W_0(k)$.
It follows that
\begin{equation} \label{eq:avg-W0}
\sum_{v \in W_0(k)} \frac{x^T M(w v) \beta}{|W_0(k)|} \ = \ c \quad \text{for all $k \ge 0$.} 
\end{equation}
Let $d > 0$ be such that $|x^T A \beta| \le d$ for all $A \in \{0,1\}^{Q \times Q}$.
Then we have
\begin{equation} \label{eq:avg-W1}
\sum_{v \in W_1(k)} \frac{|x^T M(w v) \beta|}{|W_1(k)|} \ \le \ d \quad \text{for all $k \ge 0$.} 
\end{equation}

% Fix indentation
\noindent
Let $m \ge 0$.
Define $p_1(m) := \frac{|W_1(m |u|)|}{|\Sigma|^{m |u|}}$.
We can view~$p_1(m)$ as the probability of picking a word in~$W_1(m |u|)$ when a word of length~$m |u|$ is picked uniformly at random.
We have
$p_1(m) \le \left( 1 - \frac{1}{|\Sigma|^{|u|}} \right)^m$, as in order to avoid $u$ as a factor, it has to be avoided in each of the $m$ consecutive blocks of length~$|u|$.
Thus, $\lim_{m \to \infty} p_1(m) = 0$.
We have
\begin{align*}
x^T M(w) \beta
 \ &= \ x^T M(w) \overline{A} \beta \ = \ x^T M(w) \overline{A}^{m |u|} \beta \ = \ \frac{1}{|\Sigma|^{m |u|}} \sum_{v \in \Sigma^{m |u|}} x^T M(w v) \beta \\
 \ &= \ \frac{|W_0(m |u|)|}{|\Sigma|^{m |u|}} \sum_{v \in W_0(m |u|)} \frac{x^T M(w v) \beta}{|W_0(m |u|)|} \ + \ \frac{|W_1(m |u|)|}{|\Sigma|^{m |u|}} \sum_{v \in W_1(m |u|)} \frac{x^T M(w v) \beta}{|W_1(m |u|)|} \\
 \ &= \ (1-p_1(m)) \cdot c + p_1(m) \cdot \sum_{v \in W_1(m |u|)} \frac{x^T M(w v) \beta}{|W_1(m |u|)|} \qquad (\text{by \autoref{eq:avg-W0}).}
\end{align*}
With~\autoref{eq:avg-W1} it follows that
\[
 |x^T M(w) \beta - c| \ \le \ p_1(m) ( |c| + d)\,.
\]
Since this holds for all $m \ge 0$ and $\lim_{m \to \infty} p_1(m) = 0$, we conclude that $x^T M(w) \beta = c$.
\end{proof}

% Fix indentation
\noindent
Now we prove \autoref{thm:rank-weight}.

\begin{proof}[Proof of \autoref{thm:rank-weight}]
It follows from \autoref{lem:minrank-weight} and \autoref{prop:compact} that
\[
 \alpha^T M(w) \beta \ = \ r \cdot \mcw \cdot \mrw \qquad \text{for all $w \in \Sigma^*$.}
\]
With $w = \varepsilon$ we obtain $1 = \alpha^T \beta = \alpha^T M(\varepsilon) \beta = r \cdot \mcw \cdot \mrw$, as required.
\end{proof}

\subsection{Maximal pseudo-columns}\label{subsec:pseudocolumns}

In this subsection, we define maximal pseudo-columns, which are vectors that can be seen as a relaxation of the notion of maximal columns. We show that the weight of a maximal pseudo-column is equal to the weight of a maximal column, and a maximal pseudo-column is a solution of a system of linear equations, and thus can be computed efficiently. By invoking \autoref{thm:rank-weight}, this will allow us to efficiently compute $r$.

Denote by $\MCol \subseteq \{0,1\}^Q$ the set of maximal columns.
By \autoref{thm:Cesari} (bearing in mind also \autoref{cor:ranks-equal}), the vector space spanned by all maximum columns, $\spn{\MCol}$, is at least $r$-dimensional:
\begin{prop} \label{prop:r<=MC}
We have $r \le \dim \spn{\MCol}$.
\end{prop}

% Fix indentation
\noindent
One might hypothesise that $r = \dim \spn{\MCol}$ or even that all minimum-rank matrices have the same $r$ nonzero (hence, maximum) columns.
The following example shows that neither is the case in general, not even for total DFAs.

\begin{exa} \label{ex:snap1}
Consider the total DFA with $\Sigma = \{a,b\}$ and
\begin{figure}[H]
    \centering
    \begin{subfigure}[c]{0.65\textwidth} \centering
\[
M(a) = \begin{pmatrix} 1 & 0 & 0 & 0 \\ 0 & 0 & 1 & 0 \\ 0 & 0 & 1 & 0 \\ 1 & 0 & 0 & 0 \end{pmatrix} \quad
M(b) = \begin{pmatrix} 0 & 1 & 0 & 0 \\ 0 & 1 & 0 & 0 \\ 0 & 0 & 0 & 1 \\ 0 & 0 & 0 & 1 \end{pmatrix}\]
\end{subfigure}
\begin{subfigure}[c]{0.30\textwidth} \centering
\begin{tikzpicture}[node distance = 2cm]
%,
\tikzset{every state/.style={inner sep=1pt,minimum size=1.5em}}
\node[state] (1) at (0,0) {$1$};
\node[state] (2) at (1,0) {$2$};
\node[state] (3) at (1,-1) {$3$};
\node[state] (4) at (0,-1) {$4$};
\path [-stealth, thick] (1) edge [loop,out=155,in=115,looseness=8] node[left] {$a$} (1)
(2) edge [loop,out=65,in=25,looseness=8] node[right] {$b$} (2)
(3) edge [loop,out=-25,in=-65,looseness=8] node[right] {$a$} (3)
 (4) edge [loop,out=-115,in=-155,looseness=8] node[left] {$b$} (4)
 (1) edge node[above] {$b$} (2)
(2) edge node[right] {$a$} (3)
(3) edge node[below] {$b$} (4)
(4) edge node[left] (b) {$a$} (1);
\end{tikzpicture}
\end{subfigure}
\end{figure}

% Fix indentation
\noindent
By symmetry, we have $\alpha^T = \begin{pmatrix} \frac14 & \frac14 & \frac14 & \frac14 \end{pmatrix}$.
Since no word maps states $1$ and~$3$ to the same state, we have $r=2$; i.e., $a$ and $b$ are both minimum rank.
Further, $\MCol$ consists exactly of the four nonzero columns in $M(a)$ and~$M(b)$. 
Their span $\spn{\MCol}$ is $3$-dimensional, as $\begin{pmatrix}1 & -1 & 1 & -1\end{pmatrix}$ is orthogonal to each maximum column.
Thus, $r = 2 < 3 = \dim \spn{\MCol} < 4 = |\MCol|$.
\end{exa}


% Fix indentation
\noindent
Define the vector space $U := \spn{\alpha^T M(w) - \alpha^T \mid w \in \Sigma^*}$. Intuitively, it is the set of all differences of weight distributions over the states before and after a word is applied.
Notice that for all $w_1, w_2 \in \Sigma^*$ we have $\alpha^T M(w_1) - \alpha^T M(w_2) \in U$.
Later (see the proof of \autoref{lem:U-in-NC} below) we show that $U$ is closed under post-multiplication with $M(a)$ for all $a \in \Sigma$.
Such ``forward spaces'' play an important role in weighted automata; see, e.g.,~\cite{KieferWeighted20}.
Denote the orthogonal complement of~$U$ by~$U^\bot$; i.e., $U^\bot = \{y \in \R^Q \mid \forall\,w \in \Sigma^*: \ \alpha^T M(w) y = \alpha^T y\}$. Intuitively, it is the set of vectors whose weight does not change under pre-multiplication with $M(w)$ for any $w$ (where by the weight of a vector $y$ we understand $\alpha^T y$).
Clearly, $\dim U + \dim U^\bot = |Q|$.
The following proposition follows immediately from \autoref{lem:max-column-stable}.

\begin{prop} \label{prop:MC<=V}
We have $\MCol \subseteq U^\bot$.
\end{prop}

% Fix indentation
\noindent
It follows that $\spn{\MCol}$ is a subspace of~$U^\bot$.
With \autoref{prop:r<=MC}, we have $r \le \dim \spn{\MCol} \le \dim U^\bot$.
One might hypothesise that $\spn{\MCol} = U^\bot$.
The following example  shows that this is not the case in general, not even for total DFAs.

\begin{exa}\label{ex-appendix}
Consider the DFA with $\Sigma = \{a,b,c\}$ and
\[
M(a) = \begin{pmatrix} 1 & 0 & 0 & 0 \\ 1 & 0 & 0 & 0 \\ 0 & 0 & 1 & 0 \\ 0 & 0 & 1 & 0 \end{pmatrix}, \, %\quad
M(b) = \begin{pmatrix} 0 & 1 & 0 & 0 \\ 0 & 1 & 0 & 0 \\ 0 & 0 & 0 & 1 \\ 0 & 0 & 0 & 1 \end{pmatrix}, \, %\quad
M(c) = \begin{pmatrix} 0 & 0 & 1 & 0 \\ 0 & 0 & 0 & 1 \\ 1 & 0 & 0 & 0 \\ 0 & 1 & 0 & 0 \end{pmatrix}\,.
\]
\begin{center}
\begin{tikzpicture}[node distance = 2cm]
%,
\tikzset{every state/.style={inner sep=1pt,minimum size=1.5em}}
\node[state] (1) at (0,0) {$1$};
\node[state] (2) at (1.5,0) {$2$};
\node[state] (3) at (0,-1.5) {$3$};
\node[state] (4) at (1.5,-1.5) {$4$};

\path [-stealth, thick] (1) edge [loop,out=155,in=115,looseness=8] node[left] {$a$} (1)
(2) edge [loop,out=65,in=25,looseness=8] node[right] {$b$} (2)
 (3) edge [loop,out=-115,in=-155,looseness=8] node[left] {$a$} (3)
 (4) edge [loop,out=-25,in=-65,looseness=8] node[right] {$b$} (4)
 (2) edge[bend left=15] node[below] {$a$} (1)
 (1) edge[bend left=15] node[above] {$b$} (2)
 (4) edge[bend left=15] node[below] {$a$} (3)
(3) edge[bend left=15] node[above] {$b$} (4)
 (1) edge[bend left=15] node[right] {$c$} (3)
(3) edge[bend left=15] node[left] {$c$} (1)
(2) edge[bend left=15] node[right] {$c$} (4)
(4) edge[bend left=15] node[left] (b) {$c$} (2);
\end{tikzpicture}
\end{center}

% Fix indentation
\noindent
We have $\Mer(1) = \Mer(2) = \{1,2\}$ and $\Mer(3) = \Mer(4) = \{3,4\}$.
Thus, $\MCol = \{\begin{pmatrix}1 & 1 & 0 & 0\end{pmatrix}^T,\begin{pmatrix}0 & 0 & 1 & 1\end{pmatrix}^T\}$.
Hence, $\dim \spn{\MCol} = 2$.

On the other hand, by symmetry we have $\alpha^T = \begin{pmatrix} \frac14 & \frac14 & \frac14 & \frac14 \end{pmatrix}$.
For any $w \in \Sigma^*$,
\[
\alpha^T M(w) \ = \ \begin{cases}
 \begin{pmatrix} \frac14 & \frac14 & \frac14 & \frac14 \end{pmatrix} & \text{if } w \in \{c\}^* \\
 \begin{pmatrix} \frac12 & 0 & \frac12 & 0 \end{pmatrix} & \text{if the last non-$c$ letter in $w$ is $a$} \\
 \begin{pmatrix} 0 & \frac12 & 0 & \frac12 \end{pmatrix} & \text{if the last non-$c$ letter in $w$ is $b$\,.}
\end{cases}
\]
It follows that $U = \spn{\begin{pmatrix} 1 & -1 & 1 & -1 \end{pmatrix}}$.
Thus, $\dim U^\bot = 4 - 1 = 3 > 2 = \dim \spn{\MCol}$, and $\spn{\MCol}$ is a strict subspace of~$U^\bot$.
For example, the vector $\begin{pmatrix} 1 & 0 & 0 & 1 \end{pmatrix}^T$ is in $U^\bot$ but not in $\spn{\MCol}$.
\end{exa}

% Fix indentation
\noindent
Although the dimension of~$U^\bot$ does not generally equal~$r$, the vector space~$U^\bot$ turns out useful for computing~$r$.
Recall that, by \autoref{thm:rank-weight}, we can obtain~$r$ by computing~$\mcw$ (and, symmetrically,~$\mrw$). Recall that we define for $q \in Q$
\[
 \Mer(q) \ := \ \{q' \in Q \mid \exists\,S \supseteq \{q,q'\} \text{ such that } [S] \text{ is a column}\}\,.
\]
We will need the following lemma which is easy to prove.

\begin{lem}\label{lem:max-column-body}
Let $v \in \Sigma^*$ and $q \in Q$ be such that $[v \cdot q]$ is a maximal column. Then $v \cdot q' = \emptyset$ holds for all $q' \in \Mer(q) \setminus \{q\}$.
\end{lem}

% Fix indentation
\noindent
We call a vector $y \in U^\bot$ a \emph{maximal pseudo-column} if there is $q \in Q$ with $y(q) = 1$ and $y(q') = 0$ for all $q' \not\in \Mer(q)$.
This notion, which is closely related to the ``pseudo-cuts'' from \cite{KieferW19}, can be seen as a relaxation of the notion of a maximal column: clearly, every maximal column is a maximal pseudo-column, but the converse is not true, since a maximal pseudo-column is not necessarily a vector over $\{0, 1\}$, let alone a \emph{column} in the strict sense, i.e., of the form~$[w \cdot p]$. The following lemma however shows that the weight of a maximal pseudo-column is equal to the weight of a maximal column. We will later show that computing the former can be done in \NCT.



\begin{lem} \label{lem:Cas}
Let $y$ be a maximal pseudo-column.
Then $\alpha^T y = \mcw$.
\end{lem}
\begin{proof}
Let $q \in Q$ be such that $y(q) = 1$ and $y(q') = 0$ for all $q' \not\in \Mer(q)$.
Let $w \in \Sigma^*$ be such that $[w \cdot q]$ is a maximal column.
We have
\begin{align*}
\alpha^T y \
&= \ \alpha^T M(w) y && (y \in U^\bot) \\
&= \ \sum_{q' \in Q} y(q') \alpha^T [w \cdot q'] \\
&= \ \sum_{q' \in \Mer(q)} y(q') \alpha^T [w \cdot q'] && (y(q') = 0 \text{ for } q' \not\in \Mer(q)) \\
&= \ \alpha^T [w \cdot q] + \sum_{q' \in \Mer(q) \setminus \{q\}} y(q') \alpha^T [w \cdot q'] && (y(q) = 1) \\
&= \ \alpha^T [w \cdot q] && (\text{\autoref{lem:max-column-body}}) \\
&= \ \mcw && (\text{by the choice of $w, q$}). \qedhere
\end{align*}
\end{proof}



\begin{exa} \label{ex:snap2}
We continue \autoref{ex:snap1}.
We have $U = \spn{\begin{pmatrix} 1 & -1 & 1 & -1 \end{pmatrix}}$ and $\Mer(2) = \{1,2,3\}$.
Let $y = \begin{pmatrix}4/3 & 1 & -1/3 & 0\end{pmatrix}^T$.
Then $y \in U^\bot$.
Since $y(2) = 1$ and $y(4) = 0$, vector~$y$ is a maximal pseudo-column.
Thus, by \autoref{lem:Cas}, $\mcw = \alpha^T y = \begin{pmatrix} \frac14 & \frac14 & \frac14 & \frac14 \end{pmatrix} y = \frac12$.
\end{exa}

\begin{thm}\label{thm:pseud-linsyst}
Let $\Gamma$ be a basis of $U$, and let $q \in Q$.
Then the following linear system for $y \in \mathbb{R}^Q$ has a solution, and all its solutions are maximal pseudo-columns:
\begin{align*}
 \gamma^T y \ &= \ 0 \quad \text{for all } \gamma^T \in \Gamma \\
 y(q) \ &= \ 1 \\
 y(q') \ &= \ 0 \quad \text{for all } q' \not \in \Mer(q)\,.
\end{align*}
\end{thm}
\begin{proof}
By \autoref{prop:MC<=V}, any maximal column solves the linear system.
Let $y \in \mathbb{R}^Q$ be a solution of the linear system.
The equations on the first line guarantee that $y \in U^\bot$.
Then, the equations on the second and third line guarantee that $y$~is a maximal pseudo-column.
\end{proof}




\subsection{Dealing with the non-strongly connected case}\label{subsec:non-sc}

The following lemma shows that in order to compute the minimum rank we can focus on the strongly connected case.

\begin{prop} \label{prop:unamb-sc}
Let $M : \Sigma \to \{0,1\}^{Q \times Q}$ be an unambiguous matrix monoid morphism.
Suppose that $Q_1 \cup Q_2 = Q$ is a partition of~$Q$ such that for all $w \in \Sigma^*$ it holds that $[Q_2]^T M(w) [Q_1] = 0$; i.e., for all $w \in \Sigma^*$ matrix~$M(w)$ has the block form
$
M(w) = \begin{pmatrix}
    M_1(w) & M_{12}(w) \\ 0 & M_2(w)
\end{pmatrix}\,,
$
where $M_1(w) \in \{0,1\}^{Q_1 \times Q_1}$ and $M_{12}(w) \in \{0,1\}^{Q_1 \times Q_2}$ and $M_2(w) \in \{0,1\}^{Q_2 \times Q_2}$.
%For each $\rkg \in \{\rk, \rku, \rkr\}$ we have $\rkg(M) = \rkg(M_1) + \rkg(M_2)$.
We have $\rkr(M) = \rkr(M_1) + \rkr(M_2)$ and $\rku(M) = \rku(M_1) + \rku(M_2)$. %whenever~$\rkg \in \{\rk,\rku,\rkr\}$.
\end{prop}
\begin{proof}
It is a well-known property of upper-triangular block matrices that $\rkg(M(w)) \ge \rkg(M_1(w)) + \rkg(M_2(w))$; see, e.g., \cite[Chapter 0.9.4]{HornJohnson13}.
Concerning the analogous property of~$\rku$, let $M(w) = C R$ for some $0/1$ matrices $C, R$.
Then $R$~has at least $\rku(M_2(w))$ rows whose $Q_1$-components are all~$0$, and at least $\rku(M_1(w))$ rows whose $Q_1$-components are not all~$0$.
It follows that $R$ has at least $\rku(M_1(w)) + \rku(M_2(w))$ rows.
Since the factorization $M(w) = C R$ was arbitrary, we conclude that \[\rku(M(w)) \ge \rku(M_1(w)) + \rku(M_2(w)).\]

\newcommand{\myrank}{\mathsf{rk}}
Now let $\myrank \in \{\rkg, \rku\}$ and $w_1, w_2 \in \Sigma^*$ be such that $\myrank(M_1(w_1)) = r_1$ and $\myrank(M_2(w_2)) = r_2$.
With what we have shown in the first paragraph, it suffices to show that $\myrank(M(w_1 w_2)) \le r_1 + r_2$.
For some matrices $A_1, A_2, B_1, B_2$ we have
\begin{align*}
M(w_1 w_2) \ = \ M(w_1) M(w_2) \ &= \ 
\begin{pmatrix}
    M_1(w_1) & A_1 \\ 0 & A_2
\end{pmatrix}
\begin{pmatrix}
    B_1 & B_2 \\ 0 & M_2(w_2)
\end{pmatrix} \\
&= \ 
\begin{pmatrix}
    M_1(w_1) \\ 0 
\end{pmatrix}
\begin{pmatrix}
    B_1 & B_2 
\end{pmatrix}
+
\begin{pmatrix}
    A_1 \\ A_2
\end{pmatrix}
\begin{pmatrix}
   0 & M_2(w_2)
\end{pmatrix}\,.
\end{align*}
The first summand has rank at most $\myrank \begin{pmatrix}
    M_1(w_1) \\ 0 
\end{pmatrix} = r_1$, and similarly the second summand has rank at most~$r_2$.
So we can write both ($i=1,2$) summands as $C_i R_i$, where $C_i$ has at most $r_i$ columns.
Hence, 
\[
M(w_1 w_2) \ = \ C_1 R_1 + C_2 R_2 \ = \ 
\begin{pmatrix}
    C_1 & C_2 
\end{pmatrix}
\begin{pmatrix}
    R_1 \\ R_2
\end{pmatrix}\,.
\]
It follows that $\myrank(M(w_1 w_2)) \le r_1+r_2$.
\end{proof}

% Fix indentation
\noindent
By a straightforward induction it follows from \autoref{prop:unamb-sc} that the minimum rank of an unambiguous matrix monoid is the sum of the minimum ranks of its strongly connected components (where ``incomplete'' components count as having rank~$0$).


